\section{消元与解的结构：从计算答案到看见约束}
\label{sec:03-Linear-Algebra/01-Linear-Systems-and-Matrices/02}

三个方程摆在面前：

\[
\begin{aligned}
x+2y+z&=4,\\
3x+8y+5z&=18,\\
2x+4y+2z&=8.
\end{aligned}
\]

动手之前先看一眼。第三式是第一式整体乘 2，它没带来任何新约束，只是把第一句话又喊了一遍。现在把它的右端改成 9，情况反过来：前两式已经逼着 \(2x+4y+2z=8\)，第三式却坚持等于 9，系统自相矛盾，解不存在。

写下三个方程，不等于握有三条信息。可能只有两条，也可能其中两条在打架。麻烦的是这两种病都不写在脸上——右端从 8 改到 9，字面上只动了一个数字，解集却从一条直线塌成了空集。

要的不是更熟练的代入消元，是一套能把这类情况直接暴露出来的程序。

\subsection{可逆的操作换来等价的系统}

对增广矩阵 \([A\mid b]\)，允许三种行操作：交换两行；把某一行乘上非零常数；把某一行的倍数加到另一行。它们保持解集，理由都是同一个——每一种都能原路撤销。交换可以再交换回来，乘 \(c\) 可以乘 \(1/c\) 抵消，加了多少可以再减回去。可逆意味着新系统蕴含旧系统、旧系统也蕴含新系统，两边约束的是同一批 \(x\)。

这一点值得说透：行操作换掉的是你用哪些方程去描述约束，不是约束本身。方程组的长相可以彻底改变，解集一动不动。既然如此，化简就可以放心做下去，一直做到系统足够好读为止。

\subsection{主元是新信息第一次出现的位置}

程序是这样的：从左上角出发，在尚未处理的部分里挑一个非零元素作主元，用它把同一列下方的元素全部清零，然后把注意力移到右下方的子块，重复。终点叫行阶梯形——每一行的首个非零元素都比上一行的更靠右，全零行沉到底部。

对开头那个系统走一遍。第二行减去第一行的三倍、第三行减去第一行的两倍，剩下

\[
\left(
\begin{array}{ccc|c}
1&2&1&4\\
0&2&2&6\\
0&0&0&0
\end{array}
\right).
\]

具体怎么算不重要，重要的是这张表现在把结构写在了脸上。第一列和第二列各拿到一个主元，第三列一无所获；末行整行归零，是``第三个方程冗余''的一张收据。主元的个数就是矩阵的秩 \(\rank(A)\)：真正说了新东西的行有几行。

\begin{center}
\begin{tikzpicture}[x=1cm,y=1cm,font=\small]
  \fill[blue!14] (0,0) rectangle (0.9,0.7);
  \fill[blue!14] (1.8,-0.7) rectangle (2.7,0);
  \foreach \i in {0,1,2}{
    \foreach \j in {0,1,2,3}{
      \draw[black!25] (\j*0.9,-\i*0.7) rectangle (\j*0.9+0.9,-\i*0.7+0.7);
    }
    \draw[black!25] (4.05,-\i*0.7) rectangle (4.95,-\i*0.7+0.7);
  }
  \node at (0.45,0.35) {\(p_1\)};
  \node at (1.35,0.35) {\(*\)};
  \node at (2.25,0.35) {\(*\)};
  \node at (3.15,0.35) {\(*\)};
  \node at (4.50,0.35) {\(*\)};
  \node at (0.45,-0.35) {\(0\)};
  \node at (1.35,-0.35) {\(0\)};
  \node at (2.25,-0.35) {\(p_2\)};
  \node at (3.15,-0.35) {\(*\)};
  \node at (4.50,-0.35) {\(*\)};
  \node at (0.45,-1.05) {\(0\)};
  \node at (1.35,-1.05) {\(0\)};
  \node at (2.25,-1.05) {\(0\)};
  \node at (3.15,-1.05) {\(0\)};
  \node at (4.50,-1.05) {\(0\)};
  \draw[blue!70!black,very thick] (0,0.7) -- (0,0) -- (1.8,0) -- (1.8,-0.7) -- (3.6,-0.7) -- (3.6,-1.4);
  \node[font=\footnotesize] at (0.45,-1.75) {主元};
  \node[font=\footnotesize] at (1.35,-1.75) {自由};
  \node[font=\footnotesize] at (2.25,-1.75) {主元};
  \node[font=\footnotesize] at (3.15,-1.75) {自由};
  \node[font=\footnotesize] at (4.50,-1.75) {\(b\)};
\end{tikzpicture}

\emph{阶梯线以下全是零。线每往右挪一格就跳过一列，被跳过的列拿不到主元，对应的变量可以随便取。这张图上有两个主元、两个自由变量；末行的 \(b\) 分量是零，所以系统相容——若那里是个非零数，末行就在宣称 \(0=c\)。}
\end{center}

主元列上的变量被后续方程逐层锁死，没抢到主元的列对应的变量则可以任意取值，称为自由变量。上面那个三元系统里 \(z\) 自由：令 \(z=t\)，回代得 \(y=3-t\)、\(x=-2+t\)，于是

\[
\begin{pmatrix}x\\y\\z\end{pmatrix}
=\begin{pmatrix}-2\\3\\0\end{pmatrix}
+t\begin{pmatrix}1\\-1\\1\end{pmatrix},
\qquad t\in\R .
\]

这比``有无穷多解''那句话给的信息多得多。它给的是一条具体的直线：一个落脚点，加上一个沿着走怎么都不会违反任何方程的方向。

\subsection{三种形态由主元的位置决定}

二元的图能把这件事画完。两条直线在平面上只有三种相处方式，而每一种都对应一副确定的主元结构。

\begin{center}
\begin{tikzpicture}[font=\footnotesize]
  \begin{scope}[x=0.55cm,y=0.55cm]
    \draw[black!30] (-0.2,-0.2) rectangle (4.4,4.4);
    \draw[blue!65!black,thick] (-0.1,0.6) -- (4.3,3.5);
    \draw[red!65!black,thick] (0.2,4.3) -- (4.0,0.0);
    \fill (2.156,2.087) circle (1.6pt);
  \end{scope}
  \node[align=center] at (1.15,-0.95) {两列都有主元\\解唯一};
  \begin{scope}[shift={(3.1,0)},x=0.55cm,y=0.55cm]
    \draw[black!30] (-0.2,-0.2) rectangle (4.4,4.4);
    \draw[blue!65!black,thick] (-0.1,0.4) -- (4.3,3.3);
    \draw[red!65!black,thick] (-0.1,1.5) -- (4.3,4.4);
  \end{scope}
  \node[align=center] at (4.25,-0.95) {消出 \(0=c\) 的行\\无解};
  \begin{scope}[shift={(6.2,0)},x=0.55cm,y=0.55cm]
    \draw[black!30] (-0.2,-0.2) rectangle (4.4,4.4);
    \draw[blue!65!black,thick] (-0.1,0.6) -- (4.3,3.5);
    \draw[red!65!black,thick,dashed] (-0.1,0.6) -- (4.3,3.5);
    \fill (0.8,1.19) circle (1.5pt);
    \fill (2.0,1.98) circle (1.5pt);
    \fill (3.2,2.77) circle (1.5pt);
  \end{scope}
  \node[align=center] at (7.35,-1.15) {只有一个主元\\一个自由变量\\无穷多解};
\end{tikzpicture}

\emph{中间和右边这两幅图，系数矩阵完全相同——两行成比例，只有一个主元。区别全在右端：右端跟着成比例，两条直线重合；不跟着，两条直线平行，消元会把这个不一致挤成一行 \(0=c\)。所以是主元结构加上右端的相容性，共同决定了解的形态。}
\end{center}

判据因此可以按顺序念下来。消元之后先找有没有形如 \([0\ \cdots\ 0\mid c]\)、\(c\neq0\) 的行，有就无解。没有的话，再看每一个未知量列是否都拿到了主元：都拿到了，解唯一；落空几列，就有几个自由度，解有无穷多个。

用秩说同一件事更紧凑：无解等价于 \(\rank([A\mid b])>\rank(A)\)，唯一解等价于不矛盾且 \(\rank(A)=n\)。第一句和上一节的列空间判据是同一句话——把 \(b\) 添进来抬高了秩，说明它带来了原有各列够不到的方向，正是上一节那张图右半边的``拼不出来''。消元只是把这个几何事实翻译成了一行 \(0=c\)。

要提防的直觉是：方程数等于未知量数并不保证唯一解。方阵照样可能行与行线性相关而秩不足，也照样可能矛盾。唯一性看的是每一列有没有主元，不是矩阵长得方不方。

\subsection{全部解等于一个特解加零空间}

齐次系统 \(Az=0\) 的全体解称为 \(A\) 的零空间，记作 \(\nullsp(A)\)，不少教材写成 \(\mathcal N(A)\)。它对加法和数乘封闭：两个解相加还是解，解乘上常数还是解。所以它是一条过原点的直线、一张过原点的平面，或更高维的同类东西。

拿到 \(Ax=b\) 的任意一个解 \(x_p\) 之后，其余的解一个也跑不掉。若 \(z\in\nullsp(A)\)，则 \(A(x_p+z)=b+0=b\)，得到新解；反过来若 \(x\) 也是解，则 \(A(x-x_p)=0\)，差落在零空间里。两句合起来就是

\[
\set{x:Ax=b}=x_p+\nullsp(A).
\]

非齐次解集是零空间的一次平移，方向结构完全照搬，位置挪开了，因此一般不含零向量，不构成子空间。

这个分解的价值在于把两件事分开：\(x_p\) 回答``够不够得到 \(b\)''，\(\nullsp(A)\) 回答``到了之后还能怎么晃''。晃动的自由度只跟 \(A\) 有关，与右端无关——换一个 \(b\)，解集平移到别处，形状分毫不变。线性微分方程的通解、约束优化里的可行方向、统计模型中不可辨识的参数，用的都是这一个结构。

\begin{remark}
最后那一项是数据工作者会亲身撞上的。若设计矩阵有两列近乎共线（同时放进``身高（厘米）''和``身高（米）''是最赤裸的版本），零空间就不再平凡，回归系数沿某个方向随便挪都不改变拟合值。此时报出来的单个系数没有意义，尽管预测本身照旧正常。诊断的入口是秩和主元结构，不是拟合优度。正则化之所以能救场，是因为它在零空间方向上额外补了一个偏好，把不可辨识的自由度按掉，代价是引入偏差；\hyperref[sec:13-Machine-Learning/02-Linear-Models/02]{``正则化从 Ridge 到 Lasso''} 讲的正是这笔交易。
\end{remark}

\subsection{阶梯形不唯一，主元的位置唯一}

Gauss 消元停在行阶梯形、然后回代，就够解方程了。若愿意再走一步——把每个主元归一为 1，并消去主元上方的元素——得到简化行阶梯形，解可以直接从矩阵里读出来。

差别在唯一性。同一个矩阵的简化行阶梯形是唯一的，不管你怎么选主元、按什么顺序消；行阶梯形则不唯一，取决于选取和缩放。所以``阶梯形长什么样''不是矩阵的不变量，``主元落在哪几列''才是。

还有一条容易踩空的地方：行操作保持行空间，却不保持列空间——消元会把各列改得面目全非。要取列空间的一组基，得先在阶梯形里认出主元在第几列，再回到原矩阵去取这几列。哪些子空间被行操作保住、哪些被改掉，\hyperref[sec:03-Linear-Algebra/02-Vector-Spaces/03]{``矩阵的四个基本子空间：输入、输出与丢失的信息''} 会摆到一起对照。

\subsection{小主元是一场事故}

以上都默认算术精确。到了浮点数上，主元的大小直接决定误差被放大多少倍。设想某一步的候选主元是 \(10^{-16}\)，而同列下方还躺着一个 \(1.5\)。代数上用哪个都合法；用前者，消元乘子达到 \(10^{16}\) 量级，被加到其他行上的数字会把原本有效的位数整个淹掉，最后交出的``解''可能一位都不对。

对策叫部分选主元：每一列都把绝对值最大的候选换到主元位置。于是行交换的地位变了。它不再是碰上零主元时的应急措施，而是常规动作，即便当前主元在精确算术里完全够用也照换不误。下一节把消元写成矩阵乘法时，这些交换会被单独记进一个置换矩阵，而它记下的交换里，多数是为了精度，不是为了避开零。

还得承认消元答不上来的问题。它能判定有没有解、有几个自由度，却说不出``这个解可信到什么程度''。一个系统可以既非奇异又对扰动极端敏感，消元程序一声不吭地把结果交给你，看不出任何异常。衡量这种敏感性要另起一套语言，见 \hyperref[sec:08-Numerical-Analysis/01-Floating-Point-and-Error/03]{``条件数衡量问题本身的敏感性''}。结构问题与可信度问题是两笔账，本节只结清了前一笔。
